
\medskip

\begin{center}
\begin{tabularx}{\linewidth}{|X m{4cm}|}\hline
\textbf{Rappel :}&\\
~&\\
\textbf{Orientation du lutin :}

S'orienter à 90\degres : pour se déplacer vers la droite

S'orienter à 0\degres : pour se déplacer vers le haut

S'orienter à $- 90$\degres : pour se déplacer vers la gauche

S'orienter à 180\degres: pour se déplacer vers le bas&\psset{unit=1cm}
\begin{pspicture}(-2,-2)(2,2)
\psline{<->}(-1,0)(1,0)\psline{<->}(0,-1)(0,1)
\rput(-1.5,0){$- 90$\degres}\rput(1.5,0){$90$\degres}
\rput(0,-1.5){$180$\degres}\rput(0,1.5){$0$\degres}
\end{pspicture}\\ \hline
\end{tabularx}
\end{center}
\smallskip

Le chat \includegraphics[width=0.8cm]{chat} indique la position de départ.

\medskip

\begin{enumerate}
\item ~

\parbox{0.5\linewidth}{On exécute le script 1 ci-contre.

Représenter dans l'annexe 3 le chemin parcouru par le chat.}\hfill
\parbox{0.48\linewidth}{
\begin{scratch}
\blockinit{Quand \greenflag est cliqué}
\blockpen{stylo en position d’écriture}
\blockmove{s'orienter à \ovaloperator{90}}
\blockmove{avancer de \ovaloperator{80}}
\blockrepeat{répéter \ovalnum{2} fois} 
	{
	\blockmove{tourner \turnleft{} de \ovalnum{90} degrés}
	\blockmove{avancer de \ovaloperator{80}}
	}
\end{scratch}
}
\item 
	\begin{enumerate}
		\item Indiquer sur la copie le numéro du dessin correspondant au script 2 ci-dessous.
	
\parbox{0.37\linewidth}{\begin{center}Script 2 \end{center}

\begin{scratch}
\blockinit{Quand \greenflag est cliqué}
\blockvariable{mettre \ovalvariable{pas} à \ovaloperator{\ovalvariable{80}}}
\blockpen{stylo en position d’écriture}
\blockmove{s’orienter à \ovaloperator{90}}
\blockmove{avancer de \ovaloperator{pas}}
\blockrepeat{répéter \ovalnum{2} fois} 
	{
	\blockmove{tourner \turnleft{} de \ovalnum{90} degrés}
	\blockvariable{mettre \ovalvariable{pas} à \ovaloperator{\ovalvariable{pas} - \ovalnum{20}}}
	\blockmove{avancer de \ovaloperator{pas}}
	}
\end{scratch}}\hfill \parbox{0.61\linewidth}{Le côté d'un carreau mesure 20 unités.

\bigskip

\begin{center}
\psset{unit=0.4cm}
\begin{pspicture}(20,15)
%\psgrid
\multido{\n=6+1}{9}{\psline[linewidth=0.5pt](\n,10)(\n,17)}
\multido{\n=10+1}{8}{\psline[linewidth=0.5pt](6,\n)(14,\n)}
\rput(8,11){\includegraphics[width=0.3cm]{chat}}
\rput(13,16){\large 1}
\psline[linewidth=1.25pt](8,11)(8,15)(11,15)(11,13)
\multido{\n=0+1}{9}{\psline[linewidth=0.5pt](\n,0)(\n,7)}
\multido{\n=0+1}{8}{\psline[linewidth=0.5pt](0,\n)(8,\n)}
\rput(2,1){\includegraphics[width=0.3cm]{chat}}
\rput(7,6){\large 2}
\psline[linewidth=1.25pt](2,1)(6,1)(6,4)(4,4)
\multido{\n=12+1}{9}{\psline[linewidth=0.5pt](\n,0)(\n,7)}
\multido{\n=0+1}{8}{\psline[linewidth=0.5pt](12,\n)(20,\n)}
\rput(14,1){\includegraphics[width=0.3cm]{chat}}
\rput(19,6){\large 3}
\psline[linewidth=1.25pt](14,1)(18,1)(18,5)(14,5)
\end{pspicture}
\end{center}}

		\item On souhaite modifier le script 2 pour parcourir le chemin suivant:

\bigskip

\psset{unit=0.4cm}
\begin{center}
\begin{pspicture}(8,7)
\multido{\n=0+1}{9}{\psline[linewidth=0.5pt](\n,0)(\n,7)}
\multido{\n=0+1}{8}{\psline[linewidth=0.5pt](0,\n)(8,\n)}
\psline[linewidth=1.25pt](2,1)(6,1)(6,4)(4,4)(4,3)
\rput(2,1){\includegraphics[width=0.4cm]{chat}}
\end{pspicture}
\end{center}

Quelle(s) modification(s) peut-on apporter au script 2 pour parcourir ce chemin ?

	\end{enumerate}
\end{enumerate}

\begin{center}
\textbf{ANNEXES À RENDRE AVEC LA COPIE}

\bigskip

\textbf{Annexe 3  question 1}

\vspace{0,5cm}

\parbox{0.45\linewidth}{Le côté d'un carreau mesure 20 unités.}\hfill
\parbox{0.52\linewidth}{\psset{unit=0.4cm}
\begin{pspicture}(8,7)
\psgrid[gridlabels=0pt,subgriddiv=1]
\rput{0}(2,1){\includegraphics[width=0.4cm]{chat}}
\end{pspicture}}
\end{center}


